% Source: https://artofproblemsolving.com/community/c3h1388992_circle_geometry_problem
% Stable identity: aops-topic-1388992-q01
Since $AB$ is a diameter, the inscribed angle $\angle ACB$ is a right angle (Thales' theorem).

\begin{asy}
usepackage("amsmath");
size(5.5cm);
defaultpen(fontsize(10pt));
pair A=(-1,0), B=(1,0), O=(0,0), C=dir(62), D=(C.x,0);
draw(circle(O,1),linewidth(0.9));
draw(A--B,linewidth(0.9));
draw(A--C--B,blue);
draw(C--D,red);
path sq(pair P,pair u,pair v,real s){return P+s*unit(u)--P+s*unit(u)+s*unit(v)--P+s*unit(v);}
draw(sq(D,B-D,C-D,0.07)); draw(sq(C,A-C,B-C,0.08));
dot(A);dot(B);dot(C);dot(D);dot(O);
label("$A$",A,W); label("$B$",B,E); label("$C$",C,N); label("$D$",D,S); label("$O$",O,S);
label("\mbox{}",(min(currentpicture,user=true)+max(currentpicture,user=true))/2,black); // keeps the text after the figure black when LaTeX typesets the labels inline
\end{asy}

Triangles $ADC$ and $ACB$ share the angle at $A$, and $\angle ADC=90^\circ=\angle ACB$. By AA they are similar, with $A\leftrightarrow A$, $D\leftrightarrow C$, $C\leftrightarrow B$. Hence
\[
\dfrac{AD}{AC}=\dfrac{DC}{CB},
\]
which gives $AD\cdot CB=AC\cdot CD$.
